\section{LLVM Semantics and Symbolic Execution}
\label{sec:semantics}

In our symbolic executor, an LLVM program state consists of a register map, a symbolic heap, and a path condition. Execution proceeds by interpreting LLVM IR instructions into SMT expressions. For most instructions (load, store, arithmetic, conditionals, branch, and \texttt{phi}), we use standard operational semantics over QF\_AUFNIA (the theory of quantifier-free arrays, uninterpreted functions, and non-linear integer arithmetic).

%\todo{sync with implementation: H only defines physical memory.}

A distinguishing feature of our design is that we encode all \emph{bitwise} and \emph{pointer-manipulation} instructions using \emph{integer arithmetic} rather than SMT bit-vector theories. This avoids bit-blasting and keeps constraint solving within a fragment where SMT solvers perform significantly better. Several common LLVM idioms arising in paging functions are handled with specialized encodings. For example:

\begin{enumerate}
  \item \textbf{Bit masking using \texttt{and}.}  
  Expressions of the form below are compiled into a short sequence of integer shifts and divisions that extract the middle $q$ bits, rather than expanding into per-bit tests.
  \[
      x \,\&\, (0^p 1^q 0^r)
  \]
 %\sw{SW: shall we explain why this is representative?}

  \item \textbf{Shifts via multiplication and division.}  
  The left shift 
  \[
      x \ll k
  \]
  is translated as 
  \[
      (x \cdot 2^{k}) \bmod 2^{64},
  \]
  and the logical right shift uses integer division by $2^k$.  
  When $k$ is constant, the term $2^k$ is precomputed, eliminating symbolic exponentiation.

  \item \textbf{Non-overlapping \texttt{or}.}  
  When $d_1$ and $d_2$ manipulate disjoint bit regions (a common pattern in address construction),
  \[
      d_1 \,|\, d_2
  \]
  is simplified to
  \[
      d_1 + d_2,
  \]
  entirely avoiding bit-level reasoning.
\end{enumerate}

Two additional design choices contribute to the efficiency of our symbolic executor.
First, we model only the subset of LLVM instructions that appear in paging implementations, which keeps the semantics lightweight and avoids the overhead of a full LLVM formalization. Second, all 64-bit quantities (including pointers) are treated as mathematical integers, with wrap-around enforced explicitly by $\bmod 2^{64}$. 
This eliminates the bit-blasting normally associated with SMT bit-vector reasoning.
These choices match the structural invariants of paging code, where masks, shifts, and address constructions follow regular patterns that our integer encodings capture naturally.

These encodings significantly reduce the size of generated constraints and improve solver performance on paging-related code.
The complete operational semantics and full integer-arithmetic encodings can be found in~\autoref{sec:se}.
